\section{Introduction}
\label{sec:intro}

Automatic image cropping---selecting the most aesthetically pleasing
sub-region of a photograph---is a long-standing problem in computational
photography.  Classical approaches rely on hand-crafted saliency maps or
composition rules~\cite{cgs2017}, while recent learning-based methods
train dedicated networks on human-annotated crop
datasets~\cite{gaic2019, a2rl2018, gaicd2020}.  A limitation shared by
all such methods is their dependence on task-specific training data and
the associated cost of data collection and annotation.

\baseline~\cite{cropper2025} introduced a compelling alternative:
\emph{training-free} image cropping through in-context learning (ICL)
with Vision-Language Models (VLMs).  By retrieving aesthetically
annotated examples from a training set via CLIP similarity and
presenting them as interleaved image--text demonstrations, \baseline
guides a VLM to output crop coordinates for a new query image.  An
iterative refinement loop, steered by a composite scorer (VILA
aesthetic~\cite{vila2024} + Area coverage), progressively improves
the crop over multiple rounds.  On the GAICD benchmark~\cite{gaicd2020},
\baseline achieves an IoU of \textbf{0.672} with the open-source
Mantis-8B-Idefics2~\cite{mantis2024} backbone.

However, \baseline was primarily developed and tuned for Gemini 1.5
Pro~\cite{gemini2024}, a substantially more capable (closed-source) VLM.
Its design choices---retrieval strategy, temperature setting, scorer
composition---were validated on Gemini and applied to Mantis without
adaptation.  We find that this leaves significant performance on the
table, particularly for a smaller model like Mantis that benefits from
more diverse demonstrations, broader candidate search, and less biased
scoring feedback.

We identify three specific limitations:

\begin{enumerate}
    \item \textbf{Retrieval redundancy.}  Top-$S$ CLIP retrieval often
    returns semantically near-duplicate images (\eg, ten sunset
    photographs for a sunset query), producing narrow ICL
    demonstrations that fail to expose the VLM to diverse aesthetic
    strategies (Sec.~\ref{sec:diverse_icl}).

    \item \textbf{Candidate homogeneity.}  At the recommended
    temperature of $0.05$, the VLM outputs $R$ crop candidates that
    are nearly identical (coordinate standard deviation $< 5$\,px),
    rendering the ``generate $R$, pick the best'' paradigm largely
    ineffective (Sec.~\ref{sec:multitemp}).

    \item \textbf{Scorer bias.}  The Area scorer component assigns
    $\text{score} = \text{crop\_area} / \text{image\_area}$, which
    monotonically rewards larger crops and, through the refinement
    feedback loop, drives the system toward near-full-image predictions
    (Sec.~\ref{sec:scorer_debias}).
\end{enumerate}

We address these limitations with three targeted modifications:
\emph{diverse ICL selection} via KMeans clustering of CLIP embeddings,
\emph{multi-temperature candidate generation} that unions proposals from
$T{=}0.05$ and $T{=}0.8$, and \emph{scorer debiasing} by dropping the
Area component in favor of VILA-only scoring.  Each improvement is
orthogonal and can be toggled independently.  These modifications are
particularly impactful for open-source VLMs like Mantis, which---unlike
Gemini---lack the capacity to compensate for suboptimal pipeline design
through sheer model capability.

Our contributions are:
\begin{itemize}
    \item Three complementary improvements---diverse ICL retrieval,
    multi-temperature candidate generation, and scorer
    debiasing---that independently and jointly improve cropping
    performance with open-source VLMs.
    \item Comprehensive ablations, including analysis of prompt
    sensitivity (a single instruction sentence \emph{decreases} IoU by
    0.036) and Gemini-vs-Mantis design transfer failures, providing
    practical guidance for VLM-based vision systems.
    \item State-of-the-art IoU on GAICD free-form cropping with
    Mantis-8B, improving from 0.672 to \textbf{0.762}.
    \item A diagnostic analysis revealing that \baseline's Area scorer
    causes the refinement loop to converge to near-full-image crops on
    a large fraction of test images, limiting the effectiveness of
    iterative refinement.
\end{itemize}
